\documentclass{article}
\usepackage[utf8]{inputenc}
\usepackage{xcolor}
\usepackage{hyperref}

\usepackage{rotating}

\usepackage{tabls}
\usepackage{booktabs}
\usepackage{amsmath}
\usepackage{amssymb}
\usepackage{amsthm}
\usepackage{amsfonts}
\usepackage{multicol}
\usepackage{enumitem}
\usepackage{microtype}
\usepackage{tikz}
\usepackage{pgfplots}
\usetikzlibrary{positioning}
\usepackage{multirow}

\usepackage{comment}

\usepackage{graphicx}
\usepackage{wrapfig}

\theoremstyle{plain}
\newtheorem{theorem}{Theorem}
\newtheorem*{theorem*}{Theorem}
\newtheorem{lemma}[theorem]{Lemma}
\newtheorem*{lemma*}{Lemma}
\newtheorem{proposition}[theorem]{Proposition}
\newtheorem{corollary}[theorem]{Corollary}

\theoremstyle{box}
\newtheorem{definition}[theorem]{Definition}
\newtheorem*{axiom*}{Axiom}
\newtheorem{dfn}[theorem]{Definition}
\newtheorem*{definition*}{Definition}
\newtheorem*{exercise*}{Exercise}
\newtheorem{observation}[theorem]{Observation}
\newtheorem{remark}[theorem]{Remark}
\newtheorem{example}[theorem]{Example}
\newtheorem{question}[theorem]{Question}
\newtheorem*{prep-problems}{Preparation Problems}

\newcommand{\pd}[3]{\frac{\partial^{#3} {#1}}{\partial {#2}^{#3}}}
\newcommand{\fpd}[2]{\frac{\partial {#1}}{\partial {#2}}}
\newcommand{\diff}[2]{\frac{d {#1}}{d {#2}}}


\begin{document}
In Project 3 template includes background for the data and questions we will be exploring. To complete this project you will need to complete the following tasks insuring that you meet \textbf{ALL} the specifications. If \textbf{any} of the specification are not met you will receive feedback identifying the specification(s) that still require work in order to complete the project.

\begin{itemize}
\item Focus on the Data
	\begin{itemize}
	\item Task 1a: Create a histogram of the data.
	\item Task 1b: Calculate the sample mean and the sample standard deviation for the data.
	\end{itemize}
\item Focus on the Model (a family of models)\\
Consider the probability density function, $f_1(x) = \frac{\lambda^r}{\Gamma(r)}x^{r-1}e^{-\lambda x}$ with $x \geq 0$, $r>0$, and $\lambda > 0$. This is a gamma distribution. 
	\begin{itemize}
	\item Task 2a: Verify that $f_1(x)$ could be a probability density function.
		\begin{itemize}
		\item Plot $f_1(x)$ to check that $f_1(x)$ is nonegative when $x \geq 0$.
		\item Calculate $\int_{-\infty}^{\infty} f(x) dx$ in Mathematica to check that the sum of the probabilities is 1.
		\end{itemize}
	\item Task 2b: Try different values of parameters and notice how changing the parameters of the model (distribution) changes the look of the density function. Plot the distribution for at least three choices of the parameters.  Make sure you clearly indicate which parameter values you used in plot. Describe what you learned in your exploration. 	
	\item Task 2c: Calculate the mean and standard deviation of the model (or distribution) as a function of the parameters using Mathematica.
	\end{itemize}
\item Select and Use a Model (from the Family)
	\begin{itemize}
	\item Task 3: Estimate the values of the parameters using the given data.
		\begin{itemize}
		\item Task 3a: Write down the system of equations you need to solve using the sample mean, sample standard deviation, and the mean and standard deviation of the distribution as functions of the parameters.
		\item Task 3b: Solve the system of equations both numerically and analytically. Be sure to include the work for your analytic solution (you can insert a picture).
		\end{itemize}
	\item Task 4: Simulate data using the model.
		\begin{itemize}
		\item Task 4a: Create a sample of 2500 random measurements from the model. Don't forget to look at the documentation in R to make sure you are using the right parameter values.
		\item Task 4b: Plot a histogram of your simulated data.
		\item Task 4c: Calculate the sample mean and sample standard deviation of your simulated data. How does it compare with the sample mean and sample standard deviation of the given data?
		\item Task 4d: Plot the cdf using the ecdf() function in R.
		\end{itemize}
	\item Task 5a: There is a closed form expression for the cdf of this distribution (it depends on a function you are likely not familiar with). Use Mathematica to calculate the closed form (exact formula) for the cdf. Insert an image of your Mathematica Calculation.
	\item Task 5b: Compare your the plot of the cdf using the ecdf() function and the plot of the cdf using the exact formula (closed form).
	\item Task 6: Answer the following questions using the appropriate probability calculations.
		\begin{itemize}
		\item What is the probability that the amount of explosive in a sample will be between 0.1 and 0.3? (If you calculated this probability with the pdf, check your answer using the cdf. If you calculated this probability with the cdf, check your answer using the pdf.) Insert an image of your Mathematica calculations.
		\item What is the probability that the amount of explosive in a sample will be more than 0.5? (If you calculated this probability with the pdf, check your answer using the cdf. If you calculated this probability with the cdf, check your answer using the pdf.) Insert an image of your Mathematica calculations.
		\item What is the 90$^{\text{th}}$ percentile of the distribution? Use the uniroot() function in R. (If you calculated this value using the pdf, check your answer using the closed form of the cdf. If you calculated this value with the closed form of the cdf, check your answer using the pdf.)		
		\end{itemize}
	\end{itemize}
\item A Model is an Assumption
	\begin{itemize}
	\item Task 7a: Repeat Task 2, Task 3, and Task 6 for the following functions.
		\begin{enumerate}
		\item $f_2(x) = \frac{1}{\sigma\sqrt{2\pi}}e^{-\frac{1}{2}(\frac{x-\mu}{\sigma})^2}$ with $-\infty < x < \infty$, $\mu > 0$, and $\sigma > 0$. This is a normal distribution.
		\item $f_3(x) = \lambda e^{-\lambda x}$ with $x \geq 0$ and $\lambda > 0$. This is an exponential distribution.
		\item $f_4(x) = \frac{1}{b-a}$ with $a < x < b$ and $-\infty < a < b <\infty$. This is a uniform distribution.
		\end{enumerate}
	\item Task 7b: How do the probabilities calculated in Task 6 depend on the assumption made about which family of model (distribution with unknown parameters) is used? Which distribution shape (gamma, $f_1$; normal, $f_2$; exponential, $f_3$; or uniform, $f_4$) do you think best describes the given data? Make sure to provide support for your answer using the given data, the models, and your previous calculations. 
	\end{itemize}
\item Conclusions
	\begin{itemize}
	\item Task 8: Answer the question, ``Does the soil in this area meet [a specified requirement]?" Make sure to provide support for your answer using the models and your data. 
	\end{itemize}
\item Reflection
	\begin{itemize}
	\item Task 9: Reflect on your work for this project. In 1-2 paragraphs, identify/explain 2-3 key mathematical  ideas you learned (and would like to remember) and 1-3 soft skills you needed/improved/learned while working on the project.
	\end{itemize}
\end{itemize}
 
\end{document}